% Generated by paper/scripts/export_tables.py from results/words/. Do not edit by hand.
\begin{description}\footnotesize
\item[D2] Datasets used exactly as published; derived data only via scripts in data/processed/.
\item[D3] Core study is zero-API (reuse stored judge outputs); optional small prompt test later.
\item[D4] RQ1: Between Help and Harm crisis labels; fuzzy degree k/4 humans vs k/9 LLM runs vs majority.
\item[D5] RQ2: Between Help and Harm appropriateness, 2 humans vs 3 judges $\times$ 3 runs; MentalAlign replication.
\item[D6] CounselBench: descriptive expert-uncertainty analysis only.
\item[D7] Licences: no redistribution of raw/derived data; BHH research use with citation.
\item[D8] MentalAlign null/out-of-range human scores excluded at load time, counts reported.
\item[D9] Ignore \_\_MACOSX/ archive members.
\item[D10] MentalHealthBench optional.
\item[D11] BHH alignment rule. Full rule in \texttt{docs/DECISIONS.md}.
\item[D12] BHH per-run judge scores. Full rule in \texttt{docs/DECISIONS.md}.
\item[D13] Missing LLM crisis labels (27 empty labels across 21 items in the 9 validation runs) stay missing and are never read as no\_crisis. Crisis degree = k\_crisis / n\_observed; n\_observed is stored per item (minimum 6 of 9).
\item[D14] Non-string text artefacts: answer\_text is null, with artefact\_type and artefact\_value columns. Original strings are kept unchanged.
\item[D15] RQ2 rulings: verdicts use paired-difference CIs vs S2; identical systems stay "inconclusive"; the added M3b/M5 paired differences are accepted as deviations; the 10 hypothesis comparisons are reported as exploratory (no multiplicity correction).
\item[D16] PR \#22 deviations accepted; the planned post-hoc demo sensitivity is superseded by the word-meaning pivot (D17) and not implemented.
\item[D17] Pivot (2026-11-24): main study = citizen vs LLM word models (docs/ANALYSIS\_PLAN\_WORDS.md); earlier RQ1/RQ2/CounselBench results become the motivation section; PR \#23 kept as the paper template.
\item[D18] PR \#26 rulings on the LLM elicitation harness: (1) parsing: the first number must itself lie in [0, 100], otherwise re-ask, accepted as a stricter reading of the plan, with the re-ask rate reported (item 1 superseded by D19); (2) a re-ask happens in the same conversation with the same seed, so the changed context gives a new draw without new wording; (3) \texttt{num\_predict = 64} is accepted as a technical cap, and replies that hit it are reported from the \texttt{done\_reason} column (\texttt{length}, final attempt), added before the Qwen, Gemma and Mistral runs; the original Llama rows predate that column and are covered by the D19 superseded counts, while the requeued Llama rows carry it; (4) the local Ollama server on localhost is the only network endpoint, recorded in the output, and AGENTS.md is clarified accordingly.
\item[D19] Strict bare-number parsing adopted after inspecting the first real model's raw replies and before any analysis; verbose replies such as 'Step 1 ...' had been parsed as 1. Superseded rows kept in superseded.csv.
\item[D20] Word-study analysis rulings (PR \#30, PR \#31; recorded in WP7 and confirmed by the PM in the PR \#33 review): (1) RQ2: a model x prompt x term with fewer than 30 valid answers is excluded (WP4 brief), and the width ratio is undefined (NaN, no CI) when the citizen IQR is 0 (About Even); H2b compares IQRs directly, so such a term is never counted as narrower. (2) RQ3: when all four level-word memberships are 0, the nearest-centre fallback stays the primary rule; the fallbacks at m = 0 and m = 100 are harmless (all codebooks agree) and are described; the LLM gap at about m = 33 is reported as a finding; two post-hoc sensitivities, labelled as such, are added: (a) flips counted only on items decided without a fallback, (b) the more severe of the two neighbouring words as fallback.
\item[D21] \texttt{experiments/words/describe\_respondents.py} (respondent profile of CAPphrase) accepted as a descriptive addition outside the frozen plan: aggregate counts only, no hypothesis, enters no test (PM review of PR \#33).
\end{description}
